\section{METR Time Horizon：用人类完成时间标定任务难度}

\subsection{50\% time horizon 的定义}

\videofigure{figures/time-horizon-definition.jpg}{Time horizon 定义为模型在给定成功率下能够完成的任务所对应的人类完成时长。}{00:04:26--00:06:08}

把熟练人类完成任务所需时间记为 $t$，模型成功概率记为 $P_\theta(\text{success}\mid t)$。METR 用 logistic curve 拟合成功率随任务长度的变化：

$$
P_\theta(\text{success}\mid t)=
\frac{1}{1+\exp\left[\beta\left(\log t-\log H_{50}\right)\right]}.
$$

\begin{itemize}
    \item $t$：熟练人类的完成时间；
    \item $H_{50}$：模型达到 50\% 成功率时对应的任务时长；
    \item $\beta$：成功率随任务长度下降的斜率；
    \item $P_\theta$：给定 agent 配置下的经验成功概率。
\end{itemize}

Time horizon 不是模型运行了多久，而是\textbf{它能完成的任务对人类来说需要多久}。人类时间在这里充当任务复杂度的经验标尺。

\begin{importantbox}{人类时间是难度代理，不是价值代理}
一个耗时三小时的机械整理任务可能价值很低，一个十分钟的高风险判断可能价值很高。Time horizon 衡量自主执行跨度，不能直接等价为经济影响。
\end{importantbox}

\subsection{任务套件与测量流程}

\videofigure{figures/metr-task-suite.jpg}{METR 组合秒级原子操作、分钟到小时的软件与研究任务，以及较长的 ML research task。}{00:06:08--00:07:08}

\videofigure{figures/metr-method.jpg}{评估流程包括构造多尺度任务、记录人类基线、运行 agent、估计成功率并拟合 time horizon。}{00:07:08--00:11:16}

任务覆盖软件工程、网络安全、一般推理和机器学习研究。对每个任务，先让熟练专业人员独立完成并记录时间，再让 agent 在受控环境中多次尝试，由自动测试或人工规则判定成功。

\begin{knowledgebox}{为什么用 log-time 拟合}
任务时长跨越秒、分钟、小时多个数量级。对 $\log t$ 拟合能让不同尺度的任务在曲线上更均匀，也使“能力翻倍”表现为近似线性趋势。
\end{knowledgebox}

\videofigure{figures/success-vs-human-time.jpg}{模型成功率与人类完成时间呈明显负相关，说明人类时长能解释相当一部分 agent 难度。}{00:12:22--00:14:02}

图中仍有较大残差：某些耗时长但步骤规则的任务对 agent 较容易，某些短但需要精细判断或脆弱工具操作的任务反而困难。因此 time horizon 是聚合趋势，不是单题预测器。

\subsection{本章小结}

METR 用人类时间把异质任务映射到统一难度轴，再从成功率曲线读出 $H_{50}$ 或 $H_{80}$。它适合追踪长期能力趋势，但必须保留任务类型和失败模式分析。

